% ==============================================================================
\section{Introduction}
\label{sec:intro}
% ==============================================================================

Aircraft maintenance hangars represent critical bottlenecks in airline
operations.  Every hour an aircraft spends waiting for a maintenance slot
translates directly into revenue loss and schedule disruption.  Yet the spatial
organisation of maintenance positions is rarely treated as a first-class
optimisation problem: planners typically rely on experience-based rules or
simple priority queues, leaving significant performance improvements untapped.

The core difficulty arises from the \emph{blocking structure} of the hangar.
In a typical configuration, aircraft parked in \emph{front} positions physically
obstruct the entry and exit of aircraft parked in \emph{rear} positions behind
them.  Whenever a rear-position aircraft needs to enter or leave while a
front-position aircraft is being serviced, at least one \emph{manoeuvring
movement}---towing or pushing the front aircraft out of the way---is required.
Such movements are time-consuming, require ground-crew coordination, and carry a
risk of damage.  Minimising them is therefore a hard operational concern, not
merely a secondary objective.

Formally, the \emph{aircraft hangar positioning problem} (AHPP) requires
simultaneously deciding (1) which \emph{position} in the hangar to assign
each arriving aircraft to and (2) when to \emph{schedule} its maintenance
block within that position, subject to earliest-start and target-finish
time-windows, minimum separation times between consecutive aircraft at the
same position, and blocking-arc constraints that link the schedules of
aircraft at spatially related positions.  Each aircraft is treated as a
single uninterrupted maintenance block of known total duration: the
operationally relevant quantities (makespan, delays, blocking interactions)
depend only on this aggregated block, and the internal sequence of
maintenance steps of each aircraft can be reconstructed deterministically
once its start time is fixed.

The AHPP belongs to the broad family of combined assignment and scheduling
problems.  It combines elements of the \emph{resource-constrained project
scheduling problem} (RCPSP)~\citep{Blazewicz1983}, classical
\emph{machine scheduling with blocking}~\citep{Hall1996}, and
\emph{layout-dependent sequence-dependent set-up times}~\citep{Pinedo2016}.
What distinguishes the AHPP from these antecedents is the
\emph{directed blocking-arc graph}: the interference between two aircraft is
asymmetric (front blocks rear, not vice versa), it depends on the positions
assigned, and it generates movements that are themselves part of the objective
function rather than merely set-up penalties.

The contributions of this paper are the following.  We introduce a precise
graph-based model of hangar topology built on \emph{blocking arcs} and a
formal definition of the entry and exit movements they induce
(\S\ref{sec:problem}).  On top of this model we present a compact
mixed-integer linear programming formulation in which blocking movements
are encoded through a six-binary indicator that captures both entry and
exit interactions (\S\ref{sec:milp}).  To overcome the limited scalability
of the exact formulation we propose a topology-aware GRASP that uses a
\emph{blocking-load} score per position to bias heavy aircraft away from
topologically central positions, with multi-start diversification, together
with a fixed-assignment scheduler (FAS) that re-optimises timing given a
fixed or partially fixed position assignment, and a composition strategy
that guarantees the final solution is never inferior to the topology
heuristic alone (\S\ref{sec:approaches}).  Finally, we design a systematic
benchmark with six hangar topologies---from fully independent to maximally
constrained---three slack levels, and four fleet sizes up to 30 aircraft,
comprising 120 instances that permit a controlled analysis of algorithm
behaviour (\S\ref{sec:experiments}).

The rest of the paper is organized as follows.
Section~\ref{sec:problem} formally defines the AHPP.
Section~\ref{sec:milp} presents the MILP formulation.
Section~\ref{sec:approaches} describes the heuristic solving approaches.
Section~\ref{sec:experiments} reports the computational benchmark.
Section~\ref{sec:conclusions} concludes and outlines future directions.
